% DERIVED from formal_layer.json — read-only, do not edit.
\begin{definitionv}{synth_13}[POSlate system structure]
A POSlate system over a potential-outcome system \(P\), a measurable covariate carrier \(\mathcal X\), and an outcome-support size \(K\) consists of:
\begin{itemize}
\item \textbf{(Latent space.)} The latent probability space \(\Omega\) of \(P\) is standard Borel.
\item \textbf{(Outcome support.)} The integer \(K\) satisfies \(K\ge 3\), and the ordered outcome support is \(\{0,\ldots,K-1\}\).
\item \textbf{(Selected nodes.)} Five pairwise distinct nodes of \(P\), written \(X,Z,D,S,Y\).
\item \textbf{(Value spaces.)} The node \(X\) is measurably identified with \(\mathcal X\), while \(Z\), \(D\), and \(S\) are measurably identified with \(\{0,1\}\).
\item \textbf{(Outcome node.)} The node \(Y\) is measurably identified with the finite ordered support \(\{0,\ldots,K-1\}\).
\end{itemize}
Here \(Z\) is the binary instrument, \(X\) is the covariate, \(D\) is treatment receipt, \(S\) is selection, and \(Y\) is the finite ordered outcome.
\end{definitionv}

\begin{definitionv}{synth_11}[Control potential outcome]
The symbol \(Y_0\), the potential ordered outcome under treatment zero, is the map from the sample space to the \(K\) ordered outcome levels obtained by setting treatment to zero:
\[
Y_0(\omega)=Y^{D\leftarrow 0}(\omega)\in\{0,\ldots,K-1\}.
\]
\end{definitionv}

\begin{definitionv}{synth_10}[Selection under treatment one]
The potential selection indicator under treatment one is the map \(S_1:\Omega\to\{0,1\}\) defined by
\[
S_1(\omega) := S(\omega; D\leftarrow 1).
\]
\end{definitionv}

\begin{definitionv}{synth_9}[Treatment-zero selection]
Define the potential selection indicator under treatment zero by
\[
S_0(\omega) := S^{D \leftarrow 0}(\omega).
\]
Equivalently, \(S_0\) is the selection variable evaluated counterfactually when the treatment variable is set to \(0\).
\end{definitionv}

\begin{definitionv}{synth_8}[Potential treatment under instrument one]
Let \(D_1\), the potential treatment under instrument state one, be the Boolean unit-level variable on \(\Omega\) given by the treatment value that would be realized under the intervention \(Z=1\). Equivalently, for each unit \(\omega\in\Omega\),
\[
D_1(\omega)=D_{Z\leftarrow 1}(\omega).
\]
\end{definitionv}

\begin{definitionv}{synth_7}[Potential treatment under zero]
The potential treatment under instrument value zero is the binary variable \(D_0\) defined by
\[
D_0(\omega)=D(\omega; Z=0),\qquad \omega\in\Omega.
\]
\end{definitionv}

\begin{definitionv}{synth_12}[Treatment-one outcome]
The potential ordered outcome under treatment one is the map
\[
Y_1:\Omega\to\{0,\ldots,K-1\}
\]
given by evaluating the ordered outcome counterfactual under the intervention that sets treatment receipt to \(1\).
\end{definitionv}

\begin{assumptionv}{ass:iv-independence}[Instrument independence]
The potential treatments \(D_0,D_1\), potential selection indicators \(S_0,S_1\), and potential ordered outcomes \(Y_0,Y_1\) are conditionally independent of the binary instrument \(Z\) given the covariate cell \(X\):
\[
(D_0,D_1,S_0,S_1,Y_0,Y_1)\perp Z\mid X.
\]
\end{assumptionv}

\begin{assumptionv}{ass:treatment-consistency}[Treatment consistency]
Observed treatment receipt equals the potential treatment under the realized instrument state:
\[
D=D_Z
\]
almost surely.
\end{assumptionv}

\begin{assumptionv}{ass:selection-exclusion-consistency}[Selection consistency]
Observed selection equals the potential selection indicator under the realized treatment:
\[
S=S_D
\]
almost surely.
\end{assumptionv}

\begin{assumptionv}{ass:outcome-exclusion-consistency}[Outcome consistency]
Among selected units, the observed outcome equals the potential ordered outcome under the realized treatment:
\[
Y=Y_D
\]
on \(\{S=1\}\) almost surely.
\end{assumptionv}

\begin{definitionv}{synth_17}[Observed finite data law]
The observed-data map sends each latent unit \(\omega\) to the finite observed cell
\[
O(\omega)
=
\bigl(X(\omega),Z(\omega),D(\omega),S(\omega),\widetilde Y(\omega)\bigr),
\qquad
\widetilde Y(\omega)
=
\begin{cases}
Y(\omega), & S(\omega)=1,\\
\varnothing, & S(\omega)=0.
\end{cases}
\]
The observed-data law \(P_{\mathrm{obs}}\) is the pushforward of the latent law \(P\) through this observed-data map:
\[
P_{\mathrm{obs}}=P\circ O^{-1}.
\]
\end{definitionv}

\begin{definitionv}{synth_14}[Covariate cell mass]
For each covariate cell \(x\), the cell probability is
\[
p_x := P\{\omega : X(\omega)=x\}.
\]
\end{definitionv}

\begin{definitionv}{synth_1}[Complier event definition]
The complier event \(C\) is the set of units whose potential treatment equals \(0\) when \(Z=0\) and equals \(1\) when \(Z=1\):
\[
C := \{\omega : D_0(\omega)=0 \land D_1(\omega)=1\}.
\]
\end{definitionv}

\begin{definitionv}{synth_22}[Slate-benefit partial-transport system]
A slate-benefit partial-transport system is a POSlate potential-outcome structure with finite discrete covariate support \(\mathcal X\), ordered outcome support \(\mathcal Y=\{0,1,\ldots,K-1\}\) for \(K\ge3\), binary instrument \(Z\), binary treatment \(D\), selection indicator \(S\), potential treatments \(D_0,D_1\), potential selections \(S_0,S_1\), and potential outcomes \(Y_0,Y_1\). The observed variables satisfy \(D=D_Z\), \(S=S_D\), and \(Y=Y_D\) on \(\{S=1\}\), with observed element \(O=(X,Z,D,S,SY)\) and observed law \(P_{\mathrm{obs}}=\mathcal L(O)\). The system carries a latent law \(P\) under which \((D_0,D_1,S_0,S_1,Y_0,Y_1)\perp Z\mid X\), \(D_1\ge D_0\) almost surely, and for each cell with \(P(C\mid X=x)>0\), where \(C=\{D_1>D_0\}\), a direction \(d(x)\in\{-1,+1\}\) satisfies \(P\{d(x)(S_1-S_0)\ge0\mid C,X=x\}=1\). Its cellwise observable capacities are \(\ell_i(x)\) and \(h_j(x)\), with \(m(x)=\min\{\sum_i\ell_i(x),\sum_jh_j(x)\}\), aggregate survivor-complier mass \(M=\sum_xp_xm(x)>0\), and survivor-complier outcome subcouplings constrained by the exact-mass partial-transport sets \(\Gamma_x^{\ast}\).
\end{definitionv}

\begin{definitionv}{synth_16}[Instrument propensity by cell]
For each covariate cell \(x\), the conditional instrument propensity \(\pi(x)\) is the real-valued conditional probability of \(Z=1\) in that cell:
\[
\pi(x)
=
\begin{cases}
\dfrac{\mathbb P(Z=1,\,X=x)}{\mathbb P(X=x)}, & \text{if } \mathbb P(X=x)>0,\\
0, & \text{if } \mathbb P(X=x)=0.
\end{cases}
\]
\end{definitionv}

\begin{assumptionv}{ass:instrument-overlap}[Instrument propensity overlap]
For a slate-benefit partial-transport system, the instrument-overlap condition at the instrument-overlap constant \(\varepsilon_Z\) consists of:
\begin{itemize}
\item \textbf{(Positivity.)} \(0<\varepsilon_Z\).
\item \textbf{(Half-overlap.)} \(\varepsilon_Z<1/2\).
\item \textbf{(Cellwise overlap.)} For every covariate cell \(x\in\mathcal X\) with covariate-cell probability \(p_x>0\), the conditional instrument propensity \(\pi(x)\) satisfies
\[
\varepsilon_Z\le \pi(x)\le 1-\varepsilon_Z .
\]
\end{itemize}
\end{assumptionv}

\begin{assumptionv}{ass:no-defiers}[Treatment monotonicity]
Potential treatment receipt is monotone in the instrument:
\[
D_1\ge D_0
\]
almost surely.
\end{assumptionv}

\begin{assumptionv}{ass:weak-selection-monotonicity}[Weak selection monotonicity]
For every covariate cell \(x\) with \(P(C\mid X=x)>0\), the covariate-specific direction \(d(x)\in\{-1,+1\}\) satisfies
\[
P\{d(x)(S_1-S_0)\ge 0\mid C,X=x\}=1.
\]
\end{assumptionv}

\begin{assumptionv}{ass:positive-aggregate-survivors}[Positive aggregate survivors]
The aggregate survivor-complier mass satisfies
\[
M>0.
\]
\end{assumptionv}

\begin{definitionv}{def:structural-law-class}[Structural law class \(\mathcal M\)]
The structural law class is
\[
\mathcal M
=
\{P:P\text{ satisfies }\cref{obj:ass:iv-independence,obj:ass:treatment-consistency,obj:ass:selection-exclusion-consistency,obj:ass:outcome-exclusion-consistency,obj:ass:instrument-overlap,obj:ass:no-defiers,obj:ass:weak-selection-monotonicity,obj:ass:positive-aggregate-survivors}\}.
\]
This is the finite-ordered, finite-covariate specialization of the corresponding Dong--Heiler cellwise restrictions on cells with \(P(C\mid X=x)>0\). Cells with \(P(C\mid X=x)=0\) have zero capacities and are inert for the target.
\end{definitionv}

\begin{definitionv}{def:observable-capacities}[Observable capacities \(\ell_i(x),h_j(x)\)]
For each \(x\in\mathcal X\) and \(i,j\in\mathcal Y\), let
\[
\ell_i(x)
\]
and
\[
h_j(x)
\]
be the lower-arm and upper-arm selected-complier outcome capacities obtained from the corresponding conditional instrument contrasts. Define the capacity totals, the capacity-total gap, and the exact survivor-complier mass by
\[
q_0(x)=\sum_{i\in\mathcal Y}\ell_i(x),
\qquad
q_1(x)=\sum_{j\in\mathcal Y}h_j(x),
\]
\[
\Delta q(x)=q_1(x)-q_0(x),
\qquad
m(x)=\min\{q_0(x),q_1(x)\}.
\]
\end{definitionv}

\begin{theoremv}{prop:capacity-identification}[Capacity identification]
Let \(S\) be a five-node potential-outcome slate with covariate \(X\), binary \(Z,D,S\), ordered outcome support \(\mathcal Y=\{0,\ldots,K-1\}\) with \(K\ge 3\), and standard-Borel latent probability space. Let \(\varepsilon_Z\) be the instrument-overlap constant, and let \(d:\mathcal X\to\{+1,-1\}\) be the covariate-specific weak selection-monotonicity direction. Suppose that:
\begin{itemize}
\item \textbf{(IV independence.)} \cref{obj:ass:iv-independence} holds.
\item \textbf{(Treatment consistency.)} \cref{obj:ass:treatment-consistency} holds.
\item \textbf{(Selection exclusion and consistency.)} \cref{obj:ass:selection-exclusion-consistency} holds.
\item \textbf{(Outcome exclusion and consistency.)} \cref{obj:ass:outcome-exclusion-consistency} holds.
\item \textbf{(Instrument overlap.)} \cref{obj:ass:instrument-overlap} holds with constant \(\varepsilon_Z\).
\item \textbf{(No defiers.)} \cref{obj:ass:no-defiers} holds.
\item \textbf{(Weak selection monotonicity.)} \cref{obj:ass:weak-selection-monotonicity} holds with direction \(d\).
\end{itemize}
Let \(\ell_i(x)\), \(h_j(x)\), \(\Delta q(x)\), and \(m(x)\) be the observable capacity contrasts formed from the observed-data law. Then the observable capacities are nonnegative in every covariate cell,
\[
\ell_i(x)\ge 0\quad\text{and}\quad h_j(x)\ge 0
\qquad\text{for every }x\in\mathcal X\text{ and every }i,j\in\mathcal Y,
\]
all covariate-cell probabilities satisfy \(p_x\ge 0\), and, for every supported cell \(x\) with \(p_x>0\),
\[
\ell_i(x)=P(Y_0=i,S_0=1,C\mid X=x)
\]
for every \(i\in\mathcal Y\),
\[
h_j(x)=P(Y_1=j,S_1=1,C\mid X=x)
\]
for every \(j\in\mathcal Y\),
\[
\Delta q(x)=P(S_1=1,C\mid X=x)-P(S_0=1,C\mid X=x),
\]
and
\[
m(x)=P(S_0=1,S_1=1,C\mid X=x).
\]
Moreover, for every supported \(x\), \(\Delta q(x)>0\) implies \(d(x)=+1\), \(\Delta q(x)<0\) implies \(d(x)=-1\), and \(\Delta q(x)=0\) makes both directions observationally admissible under the observed-data law.
\end{theoremv}

\begin{definitionv}{def:partial-transport-polytope}[Partial-transport polytopes \(\Gamma_x^\ast,\Gamma_x\)]
For each \(x\in\mathcal X\), define
\[
\Gamma_x^{\ast}
=
\left\{
\gamma_x\ge0:
\sum_{j\in\mathcal Y}\gamma_{ij,x}\le\ell_i(x),\
\sum_{i\in\mathcal Y}\gamma_{ij,x}\le h_j(x),\
\sum_{i,j\in\mathcal Y}\gamma_{ij,x}=m(x)
\right\}.
\]
The comparison polytopes are
\[
\Gamma_x^{+}
=
\left\{
\gamma_x\ge0:
\sum_{j\in\mathcal Y}\gamma_{ij,x}=\ell_i(x),\
\sum_{i\in\mathcal Y}\gamma_{ij,x}\le h_j(x)
\right\},
\]
\[
\Gamma_x^{-}
=
\left\{
\gamma_x\ge0:
\sum_{j\in\mathcal Y}\gamma_{ij,x}\le\ell_i(x),\
\sum_{i\in\mathcal Y}\gamma_{ij,x}=h_j(x)
\right\},
\]
and
\[
\Gamma_x^{0}
=
\left\{
\gamma_x\ge0:
\sum_{j\in\mathcal Y}\gamma_{ij,x}=\ell_i(x),\
\sum_{i\in\mathcal Y}\gamma_{ij,x}=h_j(x)
\right\}.
\]
The cellwise survivor-complier coupling polytope is
\[
\Gamma_x=\Gamma_x^{\ast}.
\]
\end{definitionv}

\begin{definitionv}{def:threshold-cuts}[Threshold cuts \(B_L(x),B_U(x)\)]
For every \(x\in\mathcal X\), define the prefix and tail sums
\[
\ell_{\le t}(x)=\sum_{i\le t}\ell_i(x),
\qquad
h_{\le t}(x)=\sum_{j\le t}h_j(x),
\]
\[
\ell_{<t}(x)=\sum_{i<t}\ell_i(x),
\qquad
h_{>t}(x)=\sum_{j>t}h_j(x).
\]
The lower and upper threshold-cut values are
\[
B_L(x)
=
\max\left\{
0,\
\max_{t\in\mathcal Y}\bigl[\ell_{\le t}(x)-h_{\le t}(x)\bigr]
+\min\{\Delta q(x),0\}
\right\}
\]
and
\[
B_U(x)
=
\min\left\{
m(x),\
\min_{t\in\mathcal Y}\bigl[\ell_{<t}(x)+h_{>t}(x)\bigr]
\right\}.
\]
These formulas apply on positive-survivor zero-gap cells and on zero-survivor cells.
\end{definitionv}

\begin{definitionv}{def:identified-interval}[Identified interval \(\Theta_I(P_{\mathrm{obs}})\)]
The endpoint map is
\[
\Psi(P_{\mathrm{obs}})
=
(\theta_L,\theta_U)
=
\left(
\frac{\sum_{x\in\mathcal X}p_xB_L(x)}{M},
\frac{\sum_{x\in\mathcal X}p_xB_U(x)}{M}
\right).
\]
The sharp identified interval for the benefit probability is
\[
\Theta_I(P_{\mathrm{obs}})=[\theta_L,\theta_U].
\]
\end{definitionv}

\begin{definitionv}{synth_4}[Threshold-cut capacity aggregates]
For each covariate cell $x\in\mathcal X$ and threshold $t\in\mathcal Y=\{0,\ldots,K-1\}$, define the lower-arm prefix, strict-prefix, and upper-tail capacities by $\ell_{\le t}(x)=\sum_{i\le t}\ell_i(x)$, $\ell_{<t}(x)=\sum_{i<t}\ell_i(x)$, and $\ell_{>t}(x)=\sum_{i>t}\ell_i(x)$. Define the corresponding higher-arm prefix and upper-tail capacities by $h_{\le t}(x)=\sum_{j\le t}h_j(x)$ and $h_{>t}(x)=\sum_{j>t}h_j(x)$.
\end{definitionv}

\begin{theoremv}{prop:tie-face-collapse}[Tie face collapse]
Let \(P\) be a five-node potential-outcome system as in \cref{obj:def:structural-law-class}, with ordered outcome support of size \(K\ge 3\). Fix an instrument-overlap constant \(\varepsilon_Z\) and a covariate-specific weak selection-monotonicity direction \(d(x)\). Assume:
\begin{itemize}
\item \textbf{(Structural restrictions.)} The law satisfies \cref{obj:ass:iv-independence,obj:ass:treatment-consistency,obj:ass:selection-exclusion-consistency,obj:ass:outcome-exclusion-consistency,obj:ass:instrument-overlap,obj:ass:no-defiers,obj:ass:weak-selection-monotonicity,obj:ass:positive-aggregate-survivors}, with instrument-overlap constant \(\varepsilon_Z\) and weak selection-monotonicity direction \(d(x)\).
\item \textbf{(Observable capacities.)} Let \(c\) be the observable capacity vector formed from the observed-data law \(P_{\mathrm{obs}}\) as in \cref{obj:def:observable-capacities}, and let its validity be the validity delivered by \cref{obj:prop:capacity-identification}.
\end{itemize}
Then, for every covariate cell \(x\) with \(p_x>0\),
\[
\bigl(\Gamma_x^{+}=\Gamma_x^{-}\ \text{and}\ \Gamma_x^{-}=\Gamma_x^{0}\bigr)
\quad\Longleftrightarrow\quad
\Delta q(x)=0 .
\]
For every such cell \(x\), if \(\Delta q(x)=0\), then
\[
P(S_0=S_1,\ C\mid X=x)=P(C\mid X=x).
\]
Under the same zero-gap condition,
\[
\Gamma_x=\Gamma_x^{+},
\qquad
\Gamma_x=\Gamma_x^{-},
\qquad
\Gamma_x=\Gamma_x^{0}.
\]
Finally, for every such cell \(x\), if \(\Delta q(x)=0\), then for every threshold \(t\),
\[
\ell_{\le t}(x)-h_{\le t}(x)=h_{>t}(x)-\ell_{>t}(x).
\]
\end{theoremv}

\begin{algorithmv}{def:threshold-flow-construction}[Threshold-flow construction \(P^L,P^U\)]
The inputs are the observable cellwise capacities \(\ell_i(x)\) and \(h_j(x)\), the capacity totals \(q_0(x)\) and \(q_1(x)\), and the threshold-cut values \(B_L(x)\) and \(B_U(x)\).

\begin{enumerate}
\item For each \(x\in\mathcal X\), compute the exact mass and the threshold sums:
\[
m(x)=\min\{q_0(x),q_1(x)\},
\]
\[
\ell_{\le t}(x)=\sum_{i\le t}\ell_i(x),
\qquad
h_{\le t}(x)=\sum_{j\le t}h_j(x),
\]
\[
\ell_{<t}(x)=\sum_{i<t}\ell_i(x),
\qquad
h_{>t}(x)=\sum_{j>t}h_j(x).
\]

\item Compute the threshold-cut values
\[
B_L(x)
=
\max\left\{
0,\
\max_{t\in\mathcal Y}\bigl[\ell_{\le t}(x)-h_{\le t}(x)\bigr]
+\min\{q_1(x)-q_0(x),0\}
\right\}
\]
and
\[
B_U(x)
=
\min\left\{
m(x),\
\min_{t\in\mathcal Y}\bigl[\ell_{<t}(x)+h_{>t}(x)\bigr]
\right\}.
\]

\item For the upper endpoint, construct a maximum subflow on the nested benefit graph \(i<j\). Extend it to total mass \(m(x)\) by allocating the residual mass on the complete bipartite graph.

\item For the lower endpoint, construct a maximum subflow on the nested nonbenefit graph \(j\le i\). Extend it to total mass \(m(x)\) by allocating the residual mass on the complete bipartite graph.

\item The resulting coupling in each cell satisfies
\[
\gamma_x\in\Gamma_x^{\ast}.
\]

\item Allocate row or column residual mass to the one-sided selected-complier stratum determined by the sign of
\[
q_1(x)-q_0(x).
\]
Place the remaining complier mass in the never-selected stratum and adjoin the unchanged never-taker and always-taker components of a baseline compatible law.

\item The construction outputs endpoint-attaining latent laws \(P^L\) and \(P^U\) with endpoint map
\[
\Psi(P_{\mathrm{obs}})
=
\left(
\frac{\sum_{x\in\mathcal X}p_xB_L(x)}{M},
\frac{\sum_{x\in\mathcal X}p_xB_U(x)}{M}
\right).
\]
\end{enumerate}
\end{algorithmv}

\begin{theoremv}{thm:full-law-endpoint-attainment}[Endpoint law attainment]
Let \(P_{\mathrm{obs}}\) be the observed-data law induced by a five-node potential-outcome slate system with covariates \(X\), binary instrument \(Z\), treatment \(D\), selection \(S\), and ordered outcome \(Y\in\{0,\ldots,K-1\}\), where \(K\ge 3\). Let \(d:\mathcal X\to\{-1,+1\}\) encode the covariate-specific weak selection-monotonicity direction, with \(d(x)=+1\) corresponding to \(S_0\le S_1\) and \(d(x)=-1\) corresponding to \(S_1\le S_0\). Let \(\varepsilon_Z\) be the instrument-overlap constant. Suppose that:
\begin{itemize}
\item \textbf{(Structural restrictions.)} The latent slate system satisfies \cref{obj:ass:iv-independence}, \cref{obj:ass:no-defiers}, and \cref{obj:ass:weak-selection-monotonicity}.
\item \textbf{(Tie-safe survivor model.)} The same law satisfies \cref{obj:ass:treatment-consistency}, \cref{obj:ass:selection-exclusion-consistency}, \cref{obj:ass:outcome-exclusion-consistency}, \cref{obj:ass:instrument-overlap} at \(\varepsilon_Z\), \cref{obj:ass:no-defiers}, and \cref{obj:ass:weak-selection-monotonicity} in direction \(d\).
\item \textbf{(Positive aggregate mass.)} For the observable capacity vector \(c\) formed from \(P_{\mathrm{obs}}\) as in \cref{obj:def:observable-capacities}, the cell probabilities \(p_x\), and the survivor-complier masses \(m(x)\),
\[
M=\sum_x p_x m(x)>0.
\]
\end{itemize}
Then there exists a baseline law compatible with \(P_{\mathrm{obs}}\) and \(c\): it induces \(P_{\mathrm{obs}}\), has observable capacities \(c\), has conditional complier mass in every cell at least \(\max\{q_0(x),q_1(x)\}\), and carries the same never-taker and always-taker components used by the latent completion. For this baseline, the cellwise threshold-flow construction in \cref{obj:def:threshold-flow-construction} returns two full laws \(P^L\) and \(P^U\) over \((Y_0,Y_1,S_0,S_1,D_0,D_1,X,Z)\). These laws are endpoint witnesses for the endpoint map \(\Psi(P_{\mathrm{obs}})=(\theta_L,\theta_U)\): each induces exactly \(P_{\mathrm{obs}}\), satisfies the structural restrictions with instrument-overlap constant \(\varepsilon_Z\) and direction \(d\), and attains
\[
P^L(Y_1>Y_0\mid S_1=S_0=1,C)=\theta_L,
\qquad
P^U(Y_1>Y_0\mid S_1=S_0=1,C)=\theta_U.
\]
Moreover, for every covariate cell \(x\), \(P^L\) realizes the lower threshold latent completion generated from the baseline at \(x\), and \(P^U\) realizes the upper threshold latent completion generated from the baseline at \(x\).
\end{theoremv}

\begin{theoremv}{thm:no-selection-reduction}[No-selection reduction]*
Let the ordered outcome support have \(K\ge 3\) levels, and let the observed capacities \(c\) be formed from the observed-data law \(P_{\mathrm{obs}}\) as in \cref{obj:prop:capacity-identification}. Suppose that:
\begin{itemize}
\item \textbf{(Structural law.)} The latent law satisfies the structural restrictions in \cref{obj:def:structural-law-class}, including \cref{obj:ass:iv-independence} and \cref{obj:ass:no-defiers}, with instrument-overlap constant \(\varepsilon_Z\) and weak selection-monotonicity direction \(d(x)\).
\item \textbf{(No selection submodel.)} For every covariate cell \(x\) with \(P(C\mid X=x)>0\),
\[
P(S_0=S_1=1\mid C,X=x)=1 .
\]
\end{itemize}
Then, for every covariate cell \(x\) with \(p_x>0\),
\[
\Delta q(x)=0
\quad\text{and}\quad
\Gamma_x=\Gamma_x^{0}.
\]
Moreover, if \(m(x)>0\), then the normalized threshold-cut endpoints satisfy
\[
\frac{B_L(x)}{m(x)}
=
\max\left\{
0,\,
\max_{t\in\mathcal Y}
\frac{\ell_{\le t}(x)-h_{\le t}(x)}{m(x)}
\right\}
\]
and
\[
\frac{B_U(x)}{m(x)}
=
\min\left\{
1,\,
\min_{t\in\mathcal Y}
\frac{\ell_{<t}(x)+h_{>t}(x)}{m(x)}
\right\} .
\]
\end{theoremv}
% CAUSALSMITH-CITED-SCOPE-BEGIN thm:no-selection-reduction
\verificationfootnotetext{\textbf{Formalization scope.} This result uses the published conclusion \citep{LuDingDasgupta2018} (Proposition 2, equation (6), page 546). The cited source proof is not formalized here: Lean verifies its use and all remaining steps. If that cited conclusion is false, the portions of this result that depend on it are not certified.}
% CAUSALSMITH-CITED-SCOPE-END thm:no-selection-reduction

\begin{definitionv}{synth_5}[Cellwise benefit-mass functional]
For a survivor-complier outcome subcoupling $\gamma_x\in\mathbb R_+^{K\times K}$ in cell $x$, define the unnormalized benefit mass by $b_x(\gamma_x)=\sum_{0\le i<j\le K-1}\gamma_{ij,x}$. Thus $b_x$ counts exactly the mass assigned to outcome pairs whose treatment-one outcome level strictly exceeds the treatment-zero outcome level.
\end{definitionv}

\begin{theoremv}{thm:sharp-exact-mass-threshold-interval}[Sharp exact-mass interval]
Let \(P_{\mathrm{obs}}\) be the observed-data law generated by a five-node potential-outcome slate with ordered outcome support of size \(K\ge 3\).  Suppose the latent law satisfies the structural restrictions in \cref{obj:def:structural-law-class} for the instrument-overlap constant \(\varepsilon_Z\) and weak-selection direction \(d(x)\).  Let \(c\) be the observable capacity vector from \cref{obj:def:observable-capacities}, identified from \(P_{\mathrm{obs}}\) as in \cref{obj:prop:capacity-identification}, and suppose that
\[
M=\sum_{x\in\mathcal X}p_xm(x)>0.
\]

Then the following statements hold.
\begin{itemize}
\item \textbf{(Sharp cellwise projection.)} For every cell \(x\) with \(p_x>0\), the set of survivor-complier couplings generated by feasible full laws with observed law \(P_{\mathrm{obs}}\) is exactly
\[
\Gamma_x^{\ast}
=
\left\{\gamma_x\ge0:
\sum_j\gamma_{ij,x}\le \ell_i(x),\ 
\sum_i\gamma_{ij,x}\le h_j(x),\ 
\sum_{i,j}\gamma_{ij,x}=m(x)
\right\}.
\]
\item \textbf{(Exact-mass branch.)} For every cell \(x\) with \(p_x>0\),
\[
\Delta q(x)>0 \Rightarrow \Gamma_x^{\ast}=\Gamma_x^{+},\qquad
\Delta q(x)<0 \Rightarrow \Gamma_x^{\ast}=\Gamma_x^{-},\qquad
\Delta q(x)=0 \Rightarrow \Gamma_x^{\ast}=\Gamma_x^{0}.
\]
\item \textbf{(Threshold endpoints.)} For every cell \(x\) with \(p_x>0\),
\[
\inf_{\gamma_x\in\Gamma_x^{\ast}} b_x(\gamma_x)=B_L(x),
\qquad
\sup_{\gamma_x\in\Gamma_x^{\ast}} b_x(\gamma_x)=B_U(x).
\]
Writing \(\gamma_x^L\) and \(\gamma_x^U\) for the lower and upper cellwise threshold-flow couplings from \cref{obj:def:threshold-flow-construction},
\[
\gamma_x^L\in\Gamma_x^{\ast},\qquad
b_x(\gamma_x^L)=B_L(x),\qquad
\gamma_x^U\in\Gamma_x^{\ast},\qquad
b_x(\gamma_x^U)=B_U(x).
\]
\item \textbf{(Zero-mass cells.)} For every cell \(x\) with \(m(x)=0\),
\[
\Gamma_x^{\ast}=\{0\}.
\]
\end{itemize}
Consequently, the sharp identified interval for the benefit probability over feasible full laws is
\[
\Theta_I(P_{\mathrm{obs}})
=
\left[
\frac{\sum_{x\in\mathcal X}p_xB_L(x)}{M},
\frac{\sum_{x\in\mathcal X}p_xB_U(x)}{M}
\right].
\]
\end{theoremv}

\begin{definitionv}{def:three-level-witness}[Three-level witness \(P_{\mathrm{obs}}^{\star}\) and \(P^{\star}\)]
The observed law \(P_{\mathrm{obs}}^{\star}\) and latent law \(P^{\star}\) are defined as follows. The covariate satisfies \(X=x^{\star}\) almost surely, the binary instrument \(Z\) satisfies
\[
P_{\mathrm{obs}}^{\star}(Z=1)=P_{\mathrm{obs}}^{\star}(Z=0)=\frac12,
\]
and the ordered outcome has \(K=3\) levels. Conditional on \(Z=0\),
\[
P_{\mathrm{obs}}^{\star}(D,S,Y) =
\begin{cases}
0.075, & (D,S,Y)=(0,1,0),\\
0.10, & (D,S,Y)=(0,1,1),\\
0.075, & (D,S,Y)=(0,1,2),\\
0.75, & (D,S,Y)=(0,0,\star),
\end{cases}
\]
with all other conditional cells assigned probability zero. Conditional on \(Z=1\),
\[
P_{\mathrm{obs}}^{\star}(D,S,Y) =
\begin{cases}
0.10, & (D,S,Y)=(1,1,0),\\
0.25, & (D,S,Y)=(1,1,1),\\
0.15, & (D,S,Y)=(1,1,2),\\
0.50, & (D,S,Y)=(1,0,\star),
\end{cases}
\]
with all other conditional cells assigned probability zero.

The latent law \(P^{\star}\) assigns every unit potential treatments \(D_0=0\) and \(D_1=1\). It allocates masses \(0.075,0.10,0.075\) to units with potential selection indicators \(S_0=S_1=1\) and diagonal potential-outcome pairs
\[
(Y_0,Y_1)=(0,0),(1,1),(2,2),
\]
respectively. It allocates mass \(0.25\) to units with \(S_0=0\) and \(S_1=1\), with residual \(Y_1\) masses
\[
(0.025,0.15,0.075)
\]
on outcome levels \(0,1,2\). It allocates mass \(0.50\) to units with \(S_0=S_1=0\). The instrument \(Z\) is independent of these potential variables with probability \(1/2\), and latent outcomes are assigned arbitrarily wherever selection makes them unobserved.
\end{definitionv}

\begin{theoremv}{prop:three-level-witness}[Three-level witness sharpness]
For the single-cell three-level witness, let \(P_{\mathrm{obs}}^{\star}\) be the observed-data law assigning mass
\[
\frac{3}{80},\frac{1}{20},\frac{3}{80},\frac{3}{8},\frac{1}{20},\frac{1}{8},\frac{3}{40},\frac{1}{4}
\]
to the eight displayed observed cells, in order.  Let \(c\) be the observable capacity vector obtained from \(P_{\mathrm{obs}}^{\star}\).  Then the following assertions hold:
\begin{itemize}
\item \textbf{(Full-law realization.)} There exists a latent potential-outcome system \(P^{\star}\) with a three-level slate system such that the survivor model is tie-safe at instrument overlap \(1/4\) with the single covariate cell retained, its observed law is \(P_{\mathrm{obs}}^{\star}\), its encoded full law is the displayed full witness law, \(D_0=0\) and \(D_1=1\) almost surely, and \(S_0\le S_1\) almost surely.
\item \textbf{(Observable capacities.)} At the single covariate cell \(x^{\star}\),
\[
(\ell_0(x^{\star}),\ell_1(x^{\star}),\ell_2(x^{\star}))
=
\left(\frac{3}{40},\frac{1}{10},\frac{3}{40}\right),
\qquad
(h_0(x^{\star}),h_1(x^{\star}),h_2(x^{\star}))
=
\left(\frac{1}{10},\frac{1}{4},\frac{3}{20}\right).
\]
\item \textbf{(Derived totals and cuts.)} The same capacity vector satisfies
\[
q_0(x^{\star})=\frac{1}{4},\qquad
m(x^{\star})=\frac{1}{4},\qquad
q_1(x^{\star})=\frac{1}{2},\qquad
\Delta q(x^{\star})=\frac{1}{4},
\]
and
\[
B_L(x^{\star})=0,
\qquad
B_U(x^{\star})=\frac{7}{40}.
\]
\item \textbf{(Identified interval.)} With the unit cell weight \(p_{x^{\star}}=1\), the nonnegative-capacity validity certificate, positive aggregate mass, and the observed-law-domain certificate for \(P_{\mathrm{obs}}^{\star}\), the sharp identified interval is
\[
\Theta_I(P_{\mathrm{obs}}^{\star})
=
\left[0,\frac{7}{10}\right].
\]
\item \textbf{(Endpoint attainment.)} There exist latent laws \(P^L\) and \(P^U\), each with a three-level slate system, satisfying the endpoint-witness conditions for \(P_{\mathrm{obs}}^{\star}\), overlap \(1/4\), and the retained single covariate cell, such that \(P^L\) attains \(0\) and \(P^U\) attains \(7/10\).
\end{itemize}
\end{theoremv}

\begin{definitionv}{synth_20}[Costed threshold-flow routines]
For finite \(\mathcal X\), \(\mathcal Y=\{0,\ldots,K-1\}\), and nonnegative capacity arrays \(\ell_i(x)\) and \(h_j(x)\), the costed sparse threshold-flow routine is the cellwise algorithm that computes \(q_0(x)\), \(q_1(x)\), \(\Delta q(x)\), \(m(x)\), the prefix and tail sums, and the threshold values \(B_L(x)\) and \(B_U(x)\); constructs an upper-attaining subflow greedily on the nested graph \(i<j\) and a lower-attaining subflow greedily on the nested graph \(j\le i\); and completes each subflow to total mass \(m(x)\) by a two-pointer allocation on the residual complete bipartite graph. Its output is \((B_L(x),B_U(x))_{x\in\mathcal X}\) together with sparse allocation traces listing exactly the positive row--column allocations produced for the lower- and upper-attaining couplings in \(\Gamma_x^{\ast}\). The actual sparse operation count is the number of arithmetic operations, comparisons, pointer advances, and trace writes executed by this sparse implementation. The charged sparse operation count assigns one charge to each prefix or tail update, threshold comparison, pointer advance, residual-capacity update, and positive-allocation trace append. The costed dense threshold-flow routine performs the same threshold-flow construction and then materializes each lower- and upper-attaining \(K\times K\) coupling matrix, writing zero entries as explicit matrix entries. Its actual dense operation count is the number of arithmetic operations, comparisons, pointer advances, residual-capacity updates, and matrix writes executed by the dense implementation. Its charged dense operation count is the sparse charge plus one charge for each materialized matrix entry.
\end{definitionv}

\begin{definitionv}{synth_21}[Sparse threshold-flow operation counts]
For a covariate cell \(x\), let \(\mathrm{Ops}_{\mathrm{sparse,actual}}(x)\) be the number of arithmetic operations and comparisons actually performed by the costed sparse threshold-flow routine of \(\cref{obj:def:threshold-flow-construction}\) in cell \(x\), including the computation of \(m(x)\), the prefix and tail sums, the threshold values \(B_L(x)\) and \(B_U(x)\), and the sparse lower- and upper-attaining couplings. Let \(\mathrm{Ops}_{\mathrm{sparse,charged}}(x)\) be the operation count assigned by the routine's charging scheme in cell \(x\), where each pass over an outcome level and each pointer advance or allocation step in the greedy nested-graph subflow and residual two-pointer completion is charged to the corresponding sparse routine step.
\end{definitionv}

\begin{theoremv}{thm:linear-sparse-threshold-flow}[Sparse threshold flow]
There exist positive integer constants \(C_{\mathrm{sparse}}\) and \(C_{\mathrm{dense}}\) such that the following holds.  Let \(\mathcal X\) be a finite covariate support, let \(K\ge 3\), and let \(\ell_y(x)\) and \(h_y(x)\), the lower- and upper-arm capacity arrays, be indexed by \(x\in\mathcal X\) and \(y\in\{0,\ldots,K-1\}\).  Suppose
\[
\ell_y(x)\ge 0,\qquad h_y(x)\ge 0
\]
for every \(x\in\mathcal X\) and every \(y\in\{0,\ldots,K-1\}\).  For the threshold-flow construction in \cref{obj:def:threshold-flow-construction}, write \(B_L(x)\) and \(B_U(x)\) for the lower and upper threshold-cut values in \cref{obj:def:threshold-cuts}.  Then:
\begin{itemize}
\item \textbf{(Sparse operation bounds.)} The actual sparse operation count
\[
\sum_{x\in\mathcal X}\mathrm{Ops}_{\mathrm{sparse,actual}}(x)
\]
and the charged sparse operation count
\[
\sum_{x\in\mathcal X}\mathrm{Ops}_{\mathrm{sparse,charged}}(x)
\]
are each bounded above by
\[
C_{\mathrm{sparse}}\,|\mathcal X|\,K .
\]

\item \textbf{(Dense operation bounds.)} The actual dense operation count
\[
\sum_{x\in\mathcal X}\mathrm{Ops}_{\mathrm{sparse,actual}}(x)+2|\mathcal X|K^2
\]
and the charged dense operation count are each bounded above by
\[
C_{\mathrm{dense}}\,|\mathcal X|\,K^2 .
\]

\item \textbf{(Magnitude-invariant charged schedules.)} For any other nonnegative capacity arrays on the same \(\mathcal X\) and the same \(K\), the charged sparse operation counts are equal and the charged dense operation counts are equal.

\item \textbf{(Cellwise output.)} For every \(x\in\mathcal X\), the costed sparse threshold-flow routine returns the threshold-cut pair \((B_L(x),B_U(x))\) and the sparse lower and upper allocation traces from \cref{obj:def:threshold-flow-construction}.  The induced lower and upper couplings \(\gamma_x^L\) and \(\gamma_x^U\) satisfy
\[
|\{(i,j):\gamma_x^L(i,j)>0\}|\le 2K-1,
\qquad
|\{(i,j):\gamma_x^U(i,j)>0\}|\le 2K-1,
\]
belong to the branch-free partial-transport polytope in \cref{obj:def:partial-transport-polytope}, and attain the threshold-cut benefit masses
\[
b_x(\gamma_x^L)=B_L(x),
\qquad
b_x(\gamma_x^U)=B_U(x).
\]
\end{itemize}
\end{theoremv}

\begin{assumptionv}{ass:uniform-aggregate-survivor-bound}[Uniform aggregate survivor bound]
The uniform aggregate lower bound \(m_{\star}\) is strictly positive and bounds the aggregate survivor-complier mass \(M\) in \cref{obj:def:identified-interval} from below:
\[
0<m_{\star}
\quad\text{and}\quad
m_{\star}\le M.
\]
\end{assumptionv}

\begin{assumptionv}{ass:iid-sampling}[I.i.d. sampling]
The observations \(O_1,\ldots,O_n\) are independent draws from the observed-data law \(P_{\mathrm{obs}}\).
\end{assumptionv}

\begin{definitionv}{def:uniform-law-class}[Uniform law class \(\mathcal P_n\)]
The uniform law class at index \(n\) is
\[
\mathcal P_n
=
\left\{
P\in\mathcal M:
P \text{ satisfies the uniform aggregate-survivor condition and the independent-sampling condition at index } n
\right\}.
\]
\end{definitionv}

\begin{definitionv}{synth_15}[Empirical cell probability]
For observations \(O_0,O_1,\ldots\) taking values in observed-data records with covariate cell component \(X\), realization \(\omega\), sample size \(n\in\mathbb N\), and covariate cell \(x\), the empirical covariate-cell probability is
\[
\widehat p_x
=
\begin{cases}
\dfrac{1}{n}\sum_{r=0}^{n-1}\mathbf 1\{X(O_r(\omega))=x\}, & n\ge 1,\\
0, & n=0.
\end{cases}
\]
\end{definitionv}

\begin{algorithmv}{def:plugin-endpoint-estimator}[Plug-in endpoint estimator \(\widehat\Psi_n\)]
Given observations \(O_1,\ldots,O_n\), a finite covariate support \(\mathcal X\), outcome support \(\mathcal Y=\{0,\ldots,K-1\}\), and screening threshold \(\eta_n\), define the projected screened plug-in endpoint estimator \(\widehat\Psi_n\) by the following steps.
\begin{enumerate}
\item For each empirical conditional probability given \((X=x,Z=z)\), use the empirical ratio when the empirical denominator is positive and use the fixed value zero when that denominator is zero. These empirical conditional probabilities form the raw capacity vector
\[
\widetilde c_n=(\widetilde\ell,\widetilde h).
\]

\item Define the nonnegative capacity cone and the projected empirical capacity vector by
\[
\mathfrak C=\prod_{x\in\mathcal X}\mathbb R_+^{2K},
\qquad
\widehat c_n
=\Pi_{\mathfrak C}(\widetilde c_n)
=\arg\min_{c\in\mathfrak C}\|c-\widetilde c_n\|_2^2.
\]
With the fixed Euclidean metric, the minimizer is unique, and \(\Pi_{\mathfrak C}\) is the coordinatewise positive-part map. Write
\[
\widehat c_n=(\widehat\ell,\widehat h).
\]

\item For each \(x\in\mathcal X\), define the empirical capacity totals, capacity-total gap, and empirical exact mass by
\[
\widehat q_0(x)=\sum_i\widehat\ell_i(x),
\qquad
\widehat q_1(x)=\sum_j\widehat h_j(x),
\]
\[
\widehat{\Delta q}(x)=\widehat q_1(x)-\widehat q_0(x),
\qquad
\widehat m(x)=\min\{\widehat q_0(x),\widehat q_1(x)\}.
\]

\item For each \(x\in\mathcal X\), define the empirical lower and upper threshold-cut values by
\[
\widehat B_L(x)=
\max\left\{
0,\ 
\max_{t\in\mathcal Y}
\left[
\sum_{i\le t}\widehat\ell_i(x)
-
\sum_{j\le t}\widehat h_j(x)
\right]
+
\min\{\widehat{\Delta q}(x),0\}
\right\}
\]
and
\[
\widehat B_U(x)=
\min\left\{
\widehat m(x),\ 
\min_{t\in\mathcal Y}
\left[
\sum_{i<t}\widehat\ell_i(x)
+
\sum_{j>t}\widehat h_j(x)
\right]
\right\}.
\]

\item Define the empirical retained-cell mass score and retained-cell indicator by
\[
\widehat r_x=\widehat p_x\widehat m(x),
\qquad
\widehat I_x=\mathbf 1\{\widehat r_x>\eta_n\}.
\]

\item Define the screened empirical aggregate mass and screened empirical endpoint numerators by
\[
\widehat M_n=\sum_x \widehat I_x\,\widehat p_x\widehat m(x),
\]
\[
\widehat N_{L,n}=\sum_x \widehat I_x\,\widehat p_x\widehat B_L(x),
\qquad
\widehat N_{U,n}=\sum_x \widehat I_x\,\widehat p_x\widehat B_U(x).
\]

\item Output the projected screened plug-in endpoint estimator
\[
\widehat\Psi_n=
\begin{cases}
\left(\widehat N_{L,n}/\widehat M_n,\ \widehat N_{U,n}/\widehat M_n\right),
& \widehat M_n>0,\\
(0,1), & \widehat M_n=0.
\end{cases}
\]
\end{enumerate}
\end{algorithmv}

\begin{definitionv}{synth_18}[Guard event collection]
For a finite covariate-cell support \(\mathcal X\) and a nonnegative integer \(K\), define \(\mathcal E=\mathcal E(\mathcal X,K)\) to be the full finite collection of guard events:
\[
\mathcal E(\mathcal X,K)=\{\text{all guard-event indices for }(\mathcal X,K)\}.
\]
\end{definitionv}

\begin{definitionv}{synth_2}[Empirical event frequency]
For observed-data records \(O_0,O_1,\ldots\), each containing a covariate cell, binary instrument, treatment, selection indicator, and optional observed outcome, the empirical frequency \(\widehat P_n\) of any observable event \(E\) at sample size \(n\) and state \(\omega\) is
\[
\widehat P_n(E;\omega)
=
\frac{1}{n}\sum_{r=0}^{n-1}\mathbf 1\{E(O_r(\omega))\}.
\]
\end{definitionv}

\begin{definitionv}{synth_3}[Multiplier process]
For observations \(O_1,\ldots,O_n\), multiplier inputs \(\xi_1,\ldots,\xi_n\), sample point \(\omega\), and an observed-data cell \(o=(x,z,d,s,y)\), with \(x\in\mathcal X\), \(z,d,s\in\{0,1\}\), and \(y\in\{0,\ldots,K-1\}\) when selected and \(y=\varnothing\) when unobserved, the multiplier process \(\mathbb G_n^\xi\) is
\[
\mathbb G_n^\xi(o)
=
n^{-1/2}
\sum_{r=1}^n
\xi_r(\omega)
\left\{
\mathbf 1\{O_r(\omega)=o\}
-
\widehat P_n(o)
\right\},
\qquad
\widehat P_n(o)
=
n^{-1}\sum_{r=1}^n \mathbf 1\{O_r(\omega)=o\}.
\]
The finite guard-event index collection consists exactly of the cell events indexed by \(x\in\mathcal X\), the arm events indexed by \((x,z)\in\mathcal X\times\{0,1\}\), and the selected-outcome events indexed by \((x,z,d,k)\in\mathcal X\times\{0,1\}\times\{0,1\}\times\{0,\ldots,K-1\}\).
\end{definitionv}

\begin{algorithmv}{def:face-aware-inference-handle}[Guarded confidence set \(\mathcal C_{1-\alpha,n}\)]
Given \(\widehat\Psi_n\), \(\eta_n\), confidence level \(1-\alpha\), instrument-overlap constant \(\varepsilon_Z\), aggregate survivor-complier lower bound \(m_{\star}\), and the finite observable event collection \(\mathcal E\), define the guarded confidence set \(\mathcal C_{1-\alpha,n}\) by the following steps.
\begin{enumerate}
\item Use the retained-cell rule
\[
\widehat I_x=\mathbf 1\{\widehat r_x>\eta_n\}.
\]

\item Define the near-active-face localization tolerance by
\[
a_n=(n\eta_n)^{-1/4}.
\]
The near-active lower and upper threshold cuts, the active face of \(\min\{\widehat{\Delta q}(x),0\}\), the active face of \(\min\{\widehat q_0(x),\widehat q_1(x)\}\), and the coordinatewise nonnegativity faces generated by \(\Pi_{\mathfrak C}\) determine a computable finite diagnostic face envelope for the multiplier process \(\mathbb G_n^\xi\) at tolerance \(a_n\).

\item Let \(\delta_n\) be the maximal empirical deviation over \(\mathcal E\):
\[
\delta_n=\max_{E\in\mathcal E}\left|\widehat P_n(E)-P_{\mathrm{obs}}(E)\right|.
\]

\item Define the union-bound empirical-deviation threshold by
\[
b_{n,\alpha}=\sqrt{\frac{|\mathcal E|}{4n\alpha}}.
\]

\item Define the capacity-and-screening deterministic error envelope and deterministic endpoint-padding guard by
\[
A_{n,\alpha}
=
\frac{64K|\mathcal X|}{\varepsilon_Z^2}b_{n,\alpha}
+
|\mathcal X|\eta_n,
\]
\[
g_{n,\alpha}
=
\min\left\{1,\frac{4A_{n,\alpha}}{m_{\star}}\right\}.
\]

\item Output the conservative guarded confidence set
\[
\mathcal C_{1-\alpha,n}
=
\left[
\max\{0,\widehat\Psi_{n,L}-g_{n,\alpha}\},\
\min\{1,\widehat\Psi_{n,U}+g_{n,\alpha}\}
\right].
\]
\end{enumerate}
\end{algorithmv}

\begin{theoremv}{thm:branch-free-pointwise-directional-limit}[Branch-free endpoint limit]
Fix a slate potential-outcome system with finite covariate support and ordered outcome support of size \(K\ge 3\).  Let \(P_{\mathrm{obs}}\) be its observed-data law, let \(O_1,\ldots,O_n\) denote the observed cells, let \(d(x)\in\{-1,+1\}\) be the weak selection-monotonicity direction, let \(\varepsilon_Z\) be the instrument-overlap constant, and let \(\eta_n\) be the screening threshold from \cref{obj:def:plugin-endpoint-estimator}.  Suppose that:
\begin{itemize}
\item \textbf{(Sampling.)} For every \(n\), the first \(n\) observations satisfy \cref{obj:ass:iid-sampling} with common law \(P_{\mathrm{obs}}\).
\item \textbf{(Overlap.)} The instrument propensity satisfies \cref{obj:ass:instrument-overlap} with constant \(\varepsilon_Z\).
\item \textbf{(Structural restrictions.)} The latent law satisfies \cref{obj:ass:iv-independence,obj:ass:treatment-consistency,obj:ass:selection-exclusion-consistency,obj:ass:outcome-exclusion-consistency,obj:ass:no-defiers,obj:ass:weak-selection-monotonicity,obj:ass:positive-aggregate-survivors}, with weak selection monotonicity in direction \(d\).
\item \textbf{(Screening rate.)} For every \(n\), \(\eta_n>0\), with \(\eta_n\to0\) and \(\sqrt n\,\eta_n\to\infty\).
\item \textbf{(Limit inputs.)} The finite multinomial central limit theorem applies to the empirical observed-cell vector, and the Hadamard directional delta method applies to tangentially directionally differentiable maps.
\end{itemize}
Let \(c\) be the observable capacity vector from \cref{obj:def:observable-capacities}.  Let
\[
M=\sum_x p_x m(x),
\]
where \(m(x)\) is the exact survivor-complier mass in cell \(x\), and let \(\mathcal X_+(P)=\{x:p_xm(x)>0\}\).  Let \(\Psi(P_{\mathrm{obs}})=(\theta_L,\theta_U)\) be the endpoint map identified by \cref{obj:prop:capacity-identification}, and let \(\widehat\Psi_n\) be the projected screened plug-in estimator from \cref{obj:def:plugin-endpoint-estimator}.  Then
\[
P\!\left(\{x:\widehat r_x>\eta_n\}=\mathcal X_+(P)\right)\to 1
\]
and
\[
\widehat\Psi_n\to_P\Psi(P_{\mathrm{obs}}).
\]
Moreover, there exist a Gaussian law \(\mathbb Z_{P_{\mathrm{obs}}}\) on observed-cell score functions and a map \(D\Psi\) such that
\[
\int z\,d\mathbb Z_{P_{\mathrm{obs}}}(z)=0
\]
and, for observed cells \(a,b\),
\[
\int z(a)z(b)\,d\mathbb Z_{P_{\mathrm{obs}}}(z)
=
\begin{cases}
P_{\mathrm{obs}}(a)\{1-P_{\mathrm{obs}}(a)\}, & a=b,\\
-P_{\mathrm{obs}}(a)P_{\mathrm{obs}}(b), & a\ne b.
\end{cases}
\]
The map \(D\Psi\) is the Hadamard directional derivative, at the observed-law mass vector and on the recovered support \(\mathcal X_+(P)\), of the reduced-support endpoint map.  With this derivative,
\[
\sqrt n\{\widehat\Psi_n-\Psi(P_{\mathrm{obs}})\}
\rightsquigarrow
D\Psi(\mathbb Z_{P_{\mathrm{obs}}}).
\]
\end{theoremv}

\begin{lemmav}{lem:weighted-capacity-mass-homogeneity}[Weighted mass homogeneity]
Let \(c=(\ell,h)\) be a capacity array as in \cref{obj:def:observable-capacities}, and let \(w=(w_x)\) be real cell weights satisfying \(w_x\ge 0\) for every covariate cell \(x\).  Define the weighted capacity array \(c^w\) cellwise by
\[
\ell_i^w(x)=w_x\ell_i(x),
\qquad
h_j^w(x)=w_xh_j(x).
\]
If \(m^w(x)\) denotes the exact survivor-complier mass computed from \(c^w\) by the formulas in \cref{obj:def:observable-capacities}, then, for every covariate cell \(x\),
\[
m^w(x)=w_x m(x).
\]
\end{lemmav}

\begin{lemmav}{lem:weighted-capacity-lower-cut-homogeneity}[Weighted lower cut homogeneity]
Let \(c=(\ell,h)\) be a capacity array, and let \(w=(w_x)_x\) be a cell weight function satisfying \(w_x\ge 0\) for every cell \(x\).  Let \(c^w\) be the cellwise weighted capacity array with lower and upper coordinates
\[
\ell_i^w(x)=w_x\ell_i(x),
\qquad
h_j^w(x)=w_xh_j(x).
\]
If \(B_L^w(x)\) denotes the lower threshold-cut value of \cref{obj:def:threshold-cuts} computed from \(c^w\), and \(B_L(x)\) denotes the same lower threshold-cut value computed from \(c\), then, for every cell \(x\),
\[
B_L^w(x)=w_x B_L(x).
\]
\end{lemmav}

\begin{lemmav}{lem:weighted-capacity-upper-cut-homogeneity}[Upper cut homogeneity]
Let \(c=(\ell,h)\) be a capacity array, and let \(w:\mathcal X\to\mathbb R\) satisfy \(w_x\ge 0\) for every \(x\in\mathcal X\).  Form the cellwise weighted capacity array \(c^w\) by
\[
\ell_i^w(x)=w_x\ell_i(x),
\qquad
h_j^w(x)=w_xh_j(x).
\]
Writing \(B_U^w\) for the upper threshold cut of \cref{obj:def:threshold-cuts} computed from \(c^w\), and \(B_U\) for the same cut computed from \(c\), one has, for every \(x\in\mathcal X\),
\[
B_U^w(x)=w_x B_U(x).
\]
\end{lemmav}

\begin{definitionv}{synth_19}[Empirical observed-data law for the \(\lambda\)-indexed sample]
For each law \(\lambda\in\Lambda\) and sample \((O_{\lambda,1},\ldots,O_{\lambda,n})\), the empirical observed-data law is \(\widehat P_{n,\lambda}=n^{-1}\sum_{r=1}^n\delta_{O_{\lambda,r}}\).
\end{definitionv}

\begin{definitionv}{synth_6}[Auxiliary multiplier process]
For each law $\lambda\in\Lambda$, let $\xi_\lambda=(\xi_{\lambda,r})_{r\ge1}$ denote auxiliary multiplier variables, conditionally independent given the observed sample and satisfying conditional mean zero and conditional variance one. They generate the diagnostic multiplier process $\mathbb G_{n,\lambda}^{\xi}=n^{-1/2}\sum_{r=1}^n\xi_{\lambda,r}\{\delta_{O_{\lambda,r}}-\widehat P_{n,\lambda}\}$ used to form the finite face diagnostic. The guarded confidence interval itself is constructed from the screened endpoint estimator and the deterministic padding $g_{n,\alpha}$.
\end{definitionv}

\begin{theoremv}{thm:uniform-deterministic-guard}[Uniform deterministic guard]
Fix a family of finite-slate potential-outcome systems indexed by \(\lambda\in\Lambda\), with nonempty finite covariate support \(\mathcal X\), ordered outcome support of size \(K\), and \(K\ge 3\).  For each \(\lambda\), let \(P_{\mathrm{obs},\lambda}\) be the observed-data law of \((X,Z,D,S,\widetilde Y)\), where \(\widetilde Y=Y\) on selected units and is missing off selection.  Let \(O_{\lambda,1},O_{\lambda,2},\ldots\) be observations taking values in the observed-data cells, let \(\xi_\lambda\) be the auxiliary process used by the guarded interval in \cref{obj:def:face-aware-inference-handle}, let \(d_\lambda(x)\) be the weak selection-monotonicity direction, and let \(\mu_\lambda\) be the sampling law.

Assume:
\begin{itemize}
\item \textbf{(Nominal level.)} The nominal error level satisfies \(0<\alpha<1\).
\item \textbf{(Overlap and mass constants.)} The instrument-overlap constant and aggregate survivor-complier mass lower bound satisfy
\[
0<\varepsilon_Z<\frac12,
\qquad
m_{\star}>0 .
\]
\item \textbf{(Screening sequence.)} The screening thresholds satisfy \(\eta_n>0\) for every \(n\),
\[
\eta_n\to 0,
\qquad
\sqrt n\,\eta_n\to\infty .
\]
\item \textbf{(Uniform law class.)} For every law \(\lambda\) and every \(n\ge 1\), the system belongs to the uniform finite-slate law class of \cref{obj:def:uniform-law-class} at index \(n\), with overlap constant \(\varepsilon_Z\), aggregate survivor-complier mass bound \(m_{\star}\), direction \(d_\lambda\), sampling law \(\mu_\lambda\), and observations \(O_{\lambda,1},\ldots,O_{\lambda,n}\).
\end{itemize}

Let \(\Theta_I(P_{\mathrm{obs},\lambda})\) be the sharp identified interval for the benefit probability from \cref{obj:def:identified-interval}, and let
\[
\Psi(P_{\mathrm{obs},\lambda})
=
(\theta_{L,\lambda},\theta_{U,\lambda})
\]
be its endpoint map.  Let \(\widehat\Psi_{n,\lambda}(\omega)\) be the projected screened plug-in endpoint estimator from \cref{obj:def:plugin-endpoint-estimator}.  Let \(\mathcal C_{1-\alpha,n,\lambda}(\omega)\) and \(g_{n,\alpha}\) be the guarded confidence interval and deterministic endpoint-padding guard from \cref{obj:def:face-aware-inference-handle}.

Then, for every \(n\ge 1\),
\[
\inf_{\lambda\in\Lambda}
\mu_\lambda\!\left\{
\omega:
\Theta_I(P_{\mathrm{obs},\lambda})
\subseteq
\mathcal C_{1-\alpha,n,\lambda}(\omega)
\right\}
\ge
1-\alpha .
\]
Moreover,
\[
\sup_{\lambda\in\Lambda}
\mu_\lambda\!\left\{
\omega:
g_{n,\alpha}
<
\max\left\{
\left|\widehat\theta_{L,n,\lambda}(\omega)-\theta_{L,\lambda}\right|,
\left|\widehat\theta_{U,n,\lambda}(\omega)-\theta_{U,\lambda}\right|
\right\}
\right\}
\longrightarrow 0,
\]
where \(\widehat\Psi_{n,\lambda}=(\widehat\theta_{L,n,\lambda},\widehat\theta_{U,n,\lambda})\), and
\[
g_{n,\alpha}
=
O\!\left(n^{-1/2}+\eta_n\right).
\]

Finally, for every deterministic sequence \(L_n>0\) such that
\[
L_n\to\infty,
\qquad
L_n=o(\sqrt n),
\]
the screening choice
\[
\eta_n=n^{-1/2}L_n
\]
satisfies, for every \(n\ge 1\),
\[
\eta_n>0,
\qquad
\eta_n\to0,
\qquad
\sqrt n\,\eta_n\to\infty,
\]
and the corresponding deterministic guard radius obeys
\[
g_{n,\alpha}
=
O\!\left(n^{-1/2}L_n\right).
\]
\end{theoremv}

\begin{assumptionv}{ass:direction-margin}[Direction margin]
For every \(x\in\mathcal X\) with \(p_x>0\) and \(m(x)>0\),
\[
|\Delta q(x)|\ge \kappa_{\sigma}.
\]
\end{assumptionv}

\begin{remarkv}{oeq:uniform-face-multiplier}[Uniform face multiplier calibration]
A natural next question is whether the near-active-face directional-multiplier handle admits a calibration such that, uniformly over direction-separated triangular arrays \(\mathcal P_n\),
\[
\liminf_{n\to\infty}\inf_{P\in\mathcal P_n}
P\{\Theta_I(P_{\mathrm{obs}})\subseteq\mathcal C_{1-\alpha,n}\}
\ge 1-\alpha,
\]
while allowing arbitrary simultaneous threshold-cut ties and covariate cells whose survivor mass approaches zero and is omitted at the unnormalized threshold \(\eta_n\).
\end{remarkv}
